\section{النهج المعتمد للنمذجة الديناميكية والفرضيات الأساسية}
\label{sec:dynamic_approach}

يعتبر بناء نموذج ديناميكي دقيق تحدياً رياضياً، حيث تتشابك المعادلات بشكل كبير مع زيادة درجات الحرية. لتبسيط الحسابات الرياضية وتخفيف العبء الحسابي على وحدة التحكم دون الإخلال بالتمثيل الفيزيائي الدقيق للنظام، تم الاستناد إلى مجموعة من الفرضيات الهندسية والفيزيائية عند بناء النموذج الديناميكي لروبوت (\LR{SCARA}) المدروس:

\begin{enumerate}
    \item \textbf{تغير كتلة الوصلة الثالثة (تأثير الحمل):}
    بما أن المهمة الأساسية للروبوت هي نقل الحمولات (\LR{Pick and Place})، فإن الكتلة الفعالة للوصلة الثالثة لا تبقى ثابتة. لذلك، تم دمج كتلة الحمل $(m_L)$ بشكل صريح مع كتلة الوصلة الثالثة الأساسية $(m_3)$ لتشكيل كتلة مركبة جديدة متغيرة تعطى بالعلاقة $(m_3^* = m_3 + m_L)$. هذا الدمج يتطلب إعادة حساب شعاع مركز الثقل $(r_{p3})$ ومصفوفة القصور الذاتي للوصلة الثالثة في كل لحظة يتغير فيها الحمل.
    
    \item \textbf{المرجعية الصفرية للجاذبية الأرضية:}
    تم اعتبار الجملة الإحداثية المرتبطة بالمفصل الأول  هي الجملة المرجعية (\LR{Reference Frame}) لحسابات الطاقة الكامنة. وبناءً عليه، فإن شعاع الجاذبية الأرضية يؤثر عمودياً نحو الأسفل وبشكل موازٍ للمحور $Z_0$، ويُعطى بالصيغة الشعاعية $g = [0, 0, -g]^T$.
    
    \item \textbf{إهمال قوى كوريوليس والقوى النابذة:}
    في العديد من التطبيقات الصناعية لروبوتات (\LR{SCARA})، وبناءً على مسارات التشغيل التي لا تتطلب سرعات زاوية عالية، يمكن افتراض أن النظام يعمل ضمن سرعات حركية منخفضة إلى متوسطة. يتيح هذا الافتراض إهمال مصفوفة قوى كوريوليس والطرد المركزي $C(q, \dot{q})$ من النموذج الديناميكي العام، مما يبسط النموذج بشكل كبير.
\end{enumerate}

بناءً على صياغة لاغرانج-أويلر والفرضيات المذكورة أعلاه، يتبسط النموذج الديناميكي العام للروبوت (المجرد من المحركات) من شكله القياسي المعقد إلى الصيغة المصفوفية المخفضة التالية:

\begin{equation}
M(q)\ddot{q} + G(q) = \tau_{robot}
\label{eq:simplified_dynamic_model}
\end{equation}

حيث:
\begin{itemize}
    \item $M(q)$: مصفوفة العطالة أو القصور الذاتي للنظام بأبعاد $3 \times 3$.
    \item $G(q)$: شعاع قوى الجاذبية الأرضية بأبعاد $3 \times 1$.
    \item $\ddot{q}$: شعاع التسارعات المفصلية بأبعاد $3 \times 1$.
    \item $\tau_{robot}$: شعاع العزوم والقوى المحركة الفعالة عند المفاصل بأبعاد $3 \times 1$.
\end{itemize}

في الأقسام التالية، سيتم تفصيل الخطوات الرياضية لحساب مصفوفة العطالة وشعاع الجاذبية بالاعتماد على الثوابت الميكانيكية للوصلات واستخدام نظريات انتقال المحاور وعزوم القصور الذاتي.